\section{Justificación}

El pronóstico de la generación eléctrica aporta información para la planeación del despacho y la integración de fuentes renovables. Los registros horarios publicados por el CENACE ofrecen un caso real para estudiar patrones temporales y evaluar cuánto cambia la precisión al ampliar el horizonte de pronóstico.

La comparación con modelos de referencia permitirá determinar si la complejidad de las redes LSTM y los transformadores se traduce en mejoras de precisión bajo un protocolo reproducible. La comparación entre el pronóstico recursivo y la salida múltiple directa permitirá estudiar cómo evoluciona el error a lo largo de las 72 horas. También se considerará el costo computacional de las arquitecturas.
